% Irodov, Part II, original problems 62--72, printed pp.70--71, PDF pp.79--80.
% k_B Boltzmann constant, R universal gas constant, N_A Avogadro constant.
\thermokineticproblem{62}{70}
现代真空泵可在室温下将压强降至
$p=4\times10^{-15}\,\mathrm{atm}$。
若剩余气体是氮气，求每立方厘米内的分子数
以及分子间平均间距。
\exsolution
对理想稀薄气体，
数密度 $n=p/(k_BT)$。
取室温 $T\simeq300\,\mathrm K$、
玻尔兹曼常数 $k_B=1.380649\times10^{-23}\,\mathrm{J/K}$，
换算
$p=4\times10^{-15}(101325)
\simeq4.05\times10^{-10}\,\mathrm{Pa}$，
所以
\[
n\simeq\frac{4.05\times10^{-10}}
{(1.38\times10^{-23})(300)}
\simeq9.8\times10^{10}\,\mathrm{m^{-3}}
=\boxed{9.8\times10^4\,\mathrm{cm^{-3}}}.
\]
若以每个分子占据的平均体积 $1/n$
定义典型间距，则
\[
\boxed{\ell\simeq n^{-1/3}
\simeq2.2\times10^{-4}\,\mathrm m=0.22\,\mathrm{mm}}.
\]
\excheck $\ell$ 是估计的平均空间尺度，
不是分子硬球直径或气体分子平均自由程。

\thermokineticproblem{63}{70}
体积 $V=5.0\,\mathrm L$ 的密闭容器
含 $m=1.4\,\mathrm g$ 氮气，
温度 $T=1800\,\mathrm K$。
其中原有氮分子的 $\eta=30\%$
已分解成两个氮原子。
求气体压强。
\exsolution
分解前氮分子总摩尔数
$n_0=m/M_{\mathrm{N_2}}$，
$M_{\mathrm{N_2}}=28\,\mathrm{g/mol}$，
故 $n_0=1.4/28=0.050\,\mathrm{mol}$。
若有分数 $\eta$ 的原分子解离，
剩余分子摩尔数 $(1-\eta)n_0$，
产生氮原子摩尔数 $2\eta n_0$。
总粒子摩尔数为
$n_{\rm particles}=(1+\eta)n_0$。
按理想气体分压相加，
\[
\boxed{p=\frac{(1+\eta)mRT}{M_{\mathrm{N_2}}V}}.
\]
代入
$V=5.0\times10^{-3}\,\mathrm{m^3}$、
$R=8.314$，得到
$\boxed{p\simeq1.9\,\mathrm{atm}}$。
\excheck 解离后粒子总数增加，
但是氮元素总质量保持不变，
不能直接只计算剩余的双原子分子。

\thermokineticproblem{64}{70}
氦和氮气的混合物在标准状况下，
总质量密度为
$\rho=0.60\,\mathrm{g/L}$。
求氦原子数密度 $n_{\rm He}$。
\exsolution
设混合物中氦原子和氮分子数密度
为 $n_{\rm He},n_{\mathrm N_2}$。
标准状况下总粒子数密度
\[
n_{\rm tot}=n_{\rm He}+n_{\mathrm N_2}
=\frac{p_0}{k_BT_0}.
\]
设每个氦原子质量 $m_{\rm He}$
和每个氮分子质量 $m_{\rm N_2}$，则
\[
\rho=m_{\rm He}n_{\rm He}
+m_{\rm N_2}n_{\mathrm N_2}
=m_{\rm N_2}n_{\rm tot}
-(m_{\rm N_2}-m_{\rm He})n_{\rm He}.
\]
解出
\[
\boxed{n_{\rm He}=
\frac{m_{\rm N_2}p_0/(k_BT_0)-\rho}
{m_{\rm N_2}-m_{\rm He}}}.
\]
也可用混合物平均摩尔质量
$M_{\rm av}=\rho RT_0/p_0$，
求氦摩尔分数
$x_{\rm He}=(28-M_{\rm av})/24$。
取标准 $T_0=273.15\,\mathrm K$、$p_0=1\,\mathrm{atm}$，
得到
\[
\boxed{n_{\rm He}\simeq1.6\times10^{19}\,\mathrm{cm^{-3}}}.
\]
\excheck 标准总分子数密度约为
$2.7\times10^{19}\,\mathrm{cm^{-3}}$；
氦原子只是其中一部分。

\thermokineticproblem{65}{70}
一束平行氮分子以
$v=400\,\mathrm{m/s}$ 撞击平面墙，
入射方向与墙法线夹角为
$\theta=30^\circ$；
束内分子数密度
$n=0.9\times10^{19}\,\mathrm{cm^{-3}}$。
假设镜面完全弹性反射，求墙所受压强。
\exsolution
单位墙面积每秒碰撞的粒子数为
$\Phi=nv\cos\theta$，
因为只有速度的法向分量决定通量。
设单个氮分子质量为 $m$。
理想弹性镜面反射使法向动量从
$mv\cos\theta$ 变为其相反数，
因此单粒子给墙冲量模为
$2mv\cos\theta$。
由动量通量得
\[
\boxed{p=2nmv^2\cos^2\theta}.
\]
氮分子质量约
$m=28u=4.65\times10^{-26}\,\mathrm{kg}$，
且 $n=9.0\times10^{24}\,\mathrm{m^{-3}}$，
代入可得
\[
\boxed{p\simeq1.0\times10^5\,\mathrm{Pa}\approx1\,\mathrm{atm}}.
\]
\excheck 若散射不再是镜面弹性反射，
法向动量变化量可能不同，
不能沿用冲量 $2mv\cos\theta$。

\thermokineticproblem{66}{70}
一种气体在标准状况下密度为
$\rho=1.3\,\mathrm{mg/cm^3}$，
其中声速为 $v=330\,\mathrm{m/s}$。
求其分子的有效自由度数 $i$。
\exsolution
在理想气体中，声速满足
\[
v^2=\gamma\frac p\rho,\qquad
\gamma=\frac{C_p}{C_V}=\frac{i+2}{i}.
\]
所以
$\gamma=\rho v^2/p$，
再由 $\gamma=1+2/i$：
\[
\boxed{i=\frac2{\rho v^2/p-1}}.
\]
密度换算为
$\rho=1.3\,\mathrm{kg/m^3}$，
取标准气压
$p=1.013\times10^5\,\mathrm{Pa}$，
代入有
$\gamma\simeq1.40$，
所以
\[
\boxed{i\simeq5}.
\]
\excheck 该结果符合振动模尚未显著激发的
刚性双原子分子（3平移+2转动自由度）。

\thermokineticproblem{67}{71}
气体的声速 $v_s$ 与气体分子均方根速率
$v_{\rm rms}$ 之比是多少？
分别求：(a) 单原子气体；(b) 刚性双原子气体。
\exsolution
理想气体声速
$v_s=\sqrt{\gamma k_BT/m}$，
分子均方根速率
$v_{\rm rms}=\sqrt{3k_BT/m}$，
两式相除得到
\[
\boxed{\frac{v_s}{v_{\rm rms}}=\sqrt{\frac\gamma3}}.
\]
对单原子气体，$\gamma=5/3$，
\[
\boxed{(v_s/v_{\rm rms})_{\rm mono}
=\frac{\sqrt5}{3}\simeq0.75}.
\]
对刚性双原子气体，$\gamma=7/5$，
\[
\boxed{(v_s/v_{\rm rms})_{\rm di}
=\sqrt{\frac7{15}}\simeq0.68}.
\]
\excheck 分子热运动速率并不是声波传播速率；
二者受同一热量纲控制，却有不同系数。

\thermokineticproblem{68}{71}
由 $N$ 个原子构成的理想气体分子，
在所有平移、转动与振动自由度均被热激发的温区，
求分子的平均能量及其中平动动能所占比例。
分别讨论空间（非线性）分子与线性分子。
\exsolution
能量均分定理规定每个\emph{二次型能量项}
平均贡献 $k_BT/2$。
对非线性分子有3平移、3转动及
$3N-6$ 个振动模；每个谐振动有动、势
各一个二次项，所以
\[
i=3+3+2(3N-6)=6N-6.
\]
因此
\[
\boxed{\overline\varepsilon_{\rm nonlinear}
=\tfrac i2k_BT=3(N-1)k_BT},
\quad
\boxed{f_{\rm trans}=\frac{3k_BT/2}{3(N-1)k_BT}
=\frac1{2(N-1)}}.
\]
线性分子有3平移、2转动及
$3N-5$ 个振动模，故
\[
i=3+2+2(3N-5)=6N-5,
\]
\[
\boxed{\overline\varepsilon_{\rm linear}
=\left(3N-\frac52\right)k_BT},\quad
\boxed{f_{\rm trans}=\frac3{6N-5}}.
\]
\excheck 对非线性分子要求 $N\ge3$；
线性分子可取 $N\ge2$，
两种分子不能套用同一转动自由度计数。

\thermokineticproblem{69}{71}
在所有平动、转动、振动自由度均已热激发的
经典理想气体中，求定容摩尔热容 $C_V$ 与
绝热指数 $\gamma$：
(a) 双原子分子；
(b) 线性 $N$ 原子分子；
(c) 空间非线性 $N$ 原子分子。
\exsolution
能量均分给每摩尔内能
$U_{\rm mol}=\frac i2RT$，
故
$C_V=iR/2$、$C_p=C_V+R$，
$\gamma=1+2/i$。
(a) 双原子分子 $i=3+2+2=7$：
\[
\boxed{C_V=\frac72R,\qquad\gamma=\frac97}.
\]
(b) 线性 $N$ 原子分子
$i=6N-5$：
\[
\boxed{C_V=\left(3N-\frac52\right)R,\qquad
\gamma=\frac{6N-3}{6N-5}}.
\]
(c) 空间非线性分子
$i=6N-6$：
\[
\boxed{C_V=3(N-1)R,\qquad
\gamma=\frac{3N-2}{3N-3}}.
\]
\excheck 该经典高温极限已计入全部振动模；
常温氧、氮气通常振动模冻结，不能把
高温双原子结果 $7R/2$ 直接用于常温。

\thermokineticproblem{70}{71}
由 $N$ 原子分子组成的理想气体在
所有自由度均激发的条件下等压膨胀。
求所吸热量中用于气体膨胀功的比例，
并给出单原子气体的这一比例。
\exsolution
等压理想气体过程中
$A=\nu R\Delta T$，
$Q=\nu C_p\Delta T$，
因此
\[
\frac AQ=\frac R{C_p}
=\frac R{(i/2+1)R}=\frac2{i+2}.
\]
非线性 $N$ 原子分子 $i=6N-6$，
故
\[
\boxed{(A/Q)_{\rm nonlinear}=\frac1{3N-2}}.
\]
线性分子 $i=6N-5$，
故
\[
\boxed{(A/Q)_{\rm linear}=\frac1{3N-3/2}}.
\]
对于单原子气体 $i=3$，
\[
\boxed{(A/Q)_{\rm mono}=\frac25}.
\]
\excheck 一般不能把单原子 $N=1$
代入“非线性分子具有3个转动自由度”
的推导；单原子需单独用 $i=3$。

\thermokineticproblem{71}{71}
某气体的质量定容和定压比热容已知：
$c_V=0.65\,\mathrm{J/(g\,K)}$、
$c_p=0.91\,\mathrm{J/(g\,K)}$。
求气体的摩尔质量 $M$ 与分子的有效自由度数 $i$。
\exsolution
由摩尔 Mayer 关系 $C_p-C_V=R$，
而质量比热与摩尔热容满足
$C_p=Mc_p,C_V=Mc_V$，
所以
\[
M(c_p-c_V)=R,
\qquad
\boxed{M=\frac R{c_p-c_V}
=\frac{8.314}{0.91-0.65}\,\mathrm{g/mol}
\simeq32\,\mathrm{g/mol}}.
\]
又
$c_p/c_V=C_p/C_V=\gamma=(i+2)/i$，
于是
\[
\boxed{i=\frac2{\gamma-1}
=\frac2{c_p/c_V-1}
=\frac{2c_V}{c_p-c_V}=5}.
\]
\excheck 摩尔质量约32 g/mol，
有效自由度5符合常温刚性双原子气体。

\thermokineticproblem{72}{71}
求气体分子的有效自由度数 $i$：
(a) 已知定压摩尔热容
$C_p=29\,\mathrm{J/(mol\,K)}$；
(b) 已知沿 $pT=\text{常数}$ 路径
的摩尔热容为
$C=29\,\mathrm{J/(mol\,K)}$。
\exsolution
(a) 理想气体
$C_p=(i/2+1)R$，所以
\[
\boxed{i=2\left(\frac{C_p}{R}-1\right)
\simeq5}.
\]
(b) 一摩尔理想气体
$V=RT/p$。沿 $pT=\text{常数}$，
$p\propto1/T$，故 $V\propto T^2$，
于是
$p\,\mathrm dV=2R\,\mathrm dT$。
第一定律
$\delta Q=C_V\,\mathrm dT+p\,\mathrm dV$
给
\[
C=C_V+2R=\left(\frac i2+2\right)R.
\]
所以
\[
\boxed{i=2\left(\frac{C}{R}-2\right)
\simeq3}.
\]
\excheck (b) 不能误把 $pT=\text{常数}$
当作等压或等温过程；该路径的
多方指数为 $n=1/2$。
